%% pc-spectre-em — spectre des ondes électromagnétiques.
%% La fréquence croît vers la gauche, la longueur d'onde vers la droite ; largeurs des bandes
%% non proportionnelles (schéma), le domaine visible est agrandi en bas.
%% Le dégradé arc-en-ciel est obtenu par une suite de fines tranches colorées (pas de `shading` :
%% le SVG resterait léger, et le rendu est identique à l'écran).
\documentclass[tikz,border=4pt]{standalone}
\usepackage{liaisons}
\begin{document}
%% --- palette de l'arc-en-ciel (violet 400 nm -> rouge 800 nm) ---------------
\definecolor{arcA}{RGB}{110,30,175}
\definecolor{arcB}{RGB}{35,60,220}
\definecolor{arcC}{RGB}{0,175,200}
\definecolor{arcD}{RGB}{25,165,60}
\definecolor{arcE}{RGB}{240,225,0}
\definecolor{arcF}{RGB}{245,140,0}
\definecolor{arcG}{RGB}{215,25,25}
%% \arcenciel{x gauche}{x droite}{y bas}{y haut}
\newcommand\arcenciel[4]{%
  \foreach \k/\ca/\cb in {0/arcA/arcB,1/arcB/arcC,2/arcC/arcD,3/arcD/arcE,4/arcE/arcF,5/arcF/arcG} {
    \foreach \j in {0,1,...,5} {
      \pgfmathsetmacro\ta{(\k*6+\j)/36}
      \pgfmathsetmacro\tb{(\k*6+\j+1)/36}
      \pgfmathsetmacro\pp{100-100*\j/6}
      \fill[\ca!\pp!\cb] ({#1+(#2-#1)*\ta},#3) rectangle ({#1+(#2-#1)*\tb+0.01},#4); }}}

\begin{tikzpicture}[x=1cm,y=1cm,
  axe/.style={-{Stealth[length=6pt]}, line width=0.8pt},
  bande/.style={draw=black!60, line width=0.4pt},
  nom/.style={font=\scriptsize, text=white, align=center},
  grad/.style={font=\scriptsize},
  zoom/.style={solideF, line width=0.6pt, dash pattern=on 1.6pt off 1.6pt}]

\def\hb{1.70}   % hauteur de la barre

%% ---------------- les bandes du spectre -----------------------------------
\foreach \xa/\xb/\co in {0/1.0/{solideF!70!black}, 1.0/2.1/{solideF},
                         2.1/3.2/{solideB!80!solideF}, 4.3/5.5/{solideA!80!black},
                         5.5/7.6/{solideD}} {
  \fill[\co] (\xa,0) rectangle (\xb,\hb); }
\arcenciel{3.2}{4.3}{0}{\hb}
\draw[bande] (0,0) rectangle (7.6,\hb);
\foreach \xa in {1.0,2.1,3.2,4.3,5.5} { \draw[bande] (\xa,0) -- (\xa,\hb); }

%% noms des domaines (dans les bandes)
\node[nom, rotate=90] at (0.50,{0.5*\hb}) {Rayons gamma};
\node[nom, rotate=90] at (1.55,{0.5*\hb}) {Rayons X};
\node[nom, rotate=90] at (2.65,{0.5*\hb}) {UV};
\node[nom, rotate=90] at (4.90,{0.5*\hb}) {IR};
\node[font=\scriptsize\bfseries, text=black] at (3.75,{0.5*\hb}) {Visible};
\node[nom] at (6.55,{0.5*\hb}) {Ondes\\radiofréquences};

%% ---------------- axe des fréquences (vers la gauche) ----------------------
\foreach \xa/\lb in {1.0/{$3{,}0\times10^{19}$ Hz}, 5.5/{300 GHz}, 7.6/{30 kHz}} {
  \draw[line width=0.5pt] (\xa,\hb) -- (\xa,\hb+0.14);
  \node[grad, anchor=south, inner sep=2pt] at (\xa,\hb+0.14) {\lb}; }
\draw[axe, solideA] (7.6,\hb+0.78) -- (0,\hb+0.78);
\node[grad, anchor=south, text=solideA] at (3.8,\hb+0.82)
  {\textbf{Fréquence $\nu$ croissante}};

%% ---------------- graduations en longueur d'onde ---------------------------
\foreach \xa in {0,1.0,2.1,3.2,4.3,5.5,7.6} { \draw[line width=0.5pt] (\xa,0) -- (\xa,-0.14); }
\foreach \xa/\lb in {0/{1 fm}, 1.0/{10 pm}, 2.1/{10 nm}, 5.5/{1 mm}, 7.6/{10 km}} {
  \node[grad, anchor=north, inner sep=2pt] at (\xa,-0.14) {\lb}; }
%% bornes du visible : décalées vers l'extérieur pour laisser passer le zoom
\node[grad, anchor=north] at (2.55,-0.50) {0,4 µm};
\node[grad, anchor=north] at (4.95,-0.50) {0,8 µm};
\draw[line width=0.4pt] (3.2,-0.16) -- (2.85,-0.50);
\draw[line width=0.4pt] (4.3,-0.16) -- (4.65,-0.50);

%% ---------------- zoom sur le domaine visible ------------------------------
\draw[zoom] (3.2,0) -- (3.2,-0.90) -- (2.20,-1.78);
\draw[zoom] (4.3,0) -- (4.3,-0.90) -- (5.60,-1.78);
\arcenciel{2.20}{5.60}{-2.36}{-1.78}
\draw[bande] (2.20,-2.36) rectangle (5.60,-1.78);
\node[grad, anchor=east] at (2.12,-2.07) {\textbf{400 nm}};
\node[grad, anchor=west] at (5.68,-2.07) {\textbf{800 nm}};
\node[grad, anchor=north, text=solideF] at (3.90,-2.42) {\textbf{VISIBLE}};

%% ---------------- axe des longueurs d'onde (vers la droite) ----------------
\draw[axe] (0,-2.92) -- (7.6,-2.92);
\node[grad, anchor=north] at (3.8,-2.98) {\textbf{Longueur d'onde $\lambda$ croissante}};
\end{tikzpicture}
\end{document}
